\begin{helper}
Let \(s,\ell,m,n\) be natural numbers with \(m\le n\).  If
\(\noderef{MulticolorRamseyProperty}(s,\ell,m)\) holds, then
\(\noderef{MulticolorRamseyProperty}(s,\ell,n)\) holds.
\end{helper}
\begin{proof}
Assume \(\noderef{MulticolorRamseyProperty}(s,\ell,m)\), and let
\(\chi:E(K_n)\to[\ell]\) be any \(\ell\)-coloring of the edges on the vertex
set \([n]\).  Embed \([m]\) into \([n]\) by the initial inclusion.  Restrict
\(\chi\) along this embedding: the color of the unordered pair \(\{x,y\}\) in
\([m]\) is defined to be the color of the unordered pair of its images in
\([n]\).

By the Ramsey property at \(m\), this restricted coloring has an \(s\)-element
set \(S\subseteq[m]\) whose edges all have one color \(c\in[\ell]\).  Map
\(S\) into \([n]\) by the same inclusion.  Since the inclusion is injective,
the image still has cardinality \(s\).  If two distinct image vertices come
from \(x,y\in S\), then their edge color in the original coloring is exactly
the edge color of \(\{x,y\}\) in the restricted coloring, hence is \(c\).
Thus the image of \(S\) is a monochromatic \(s\)-clique in the original
coloring.  Because the original coloring was arbitrary,
\(\noderef{MulticolorRamseyProperty}(s,\ell,n)\) holds.
\end{proof}

